% GENERATED FILE. Edit canonical lessons, then run npm run generate:books.
% canonical-source-hash: fnv1a64:02f38eb0545b0482
% canonical-lessons: TE-C207-callarcu, TE-C207-tattu, TE-C207-tempu, TE-C207-dunnu, TE-C207-pandincu, TE-R207-first-pass-words-from-chapters-199-202, TE-R207-second-pass-words-from-chapters-203-207

\chapter{To cool, To tap, To pluck, To till, To grow}
\label{ch:te-to-cool-to-tap-to-pluck-to-till-to-grow}
\hlchaptermodality{207}

\begin{chapteropening}
\textbf{By the end of this chapter:} \emph{I can say to cool, to tap, to pluck, to till and to grow.}

Complete the last lesson of chapter 207: I can say to cool, to tap, to pluck, to till and to grow.
\end{chapteropening}

\section[\texorpdfstring{\te{చల్లార్చు} (callārcu)}{చల్లార్చు (callārcu)}]{\te{చల్లార్చు} (callārcu) — to cool (something) down}

\label{lesson:TE-C207-callarcu}



\begin{quote}
Before the new one: say the Telugu for to tread on, then the Telugu for to hit.
\end{quote}

\subsection*{You'll want to know: \te{చల్లార్చు}}
\textbf{\te{చల్లార్చు}} — \emph{callārcu} — \textquotedblleft{}to cool (something) down\textquotedblright{}.

\begin{grammarlens}[title={What you've built}]
One more word for this chapter.
\end{grammarlens}

\subsection*{Your turn}
\begin{itemize}
\raggedright
  \item \emph{Say it:} \emph{callārcu}
  \item \emph{Say it:} \emph{callārcu}, once more
  \item \emph{Say it:} say \emph{callārcu}
  \item \emph{Recall:} say the Telugu for to tread on, then the Telugu for to hit, then say \emph{callārcu} again
\end{itemize}

\begin{tcolorbox}[breakable,colback=teal!4,colframe=teal!35!black,title={Before you move on}]
What does \te{చల్లార్చు} mean? (\textquotedblleft{}To cool (something) down\textquotedblright{}.) Say it once more.
\end{tcolorbox}
\section[\texorpdfstring{\te{తట్టు} (taṭṭu)}{తట్టు (taṭṭu)}]{\te{తట్టు} (taṭṭu) — to tap, to pat}

\label{lesson:TE-C207-tattu}



\begin{quote}
Before the new one: say the Telugu for to hit, then the Telugu for to cool.
\end{quote}

\subsection*{You'll want to know: \te{తట్టు}}
\textbf{\te{తట్టు}} — \emph{taṭṭu} — \textquotedblleft{}to tap, to pat\textquotedblright{}.

\begin{grammarlens}[title={What you've built}]
One more word for this chapter.
\end{grammarlens}

\subsection*{Your turn}
\begin{itemize}
\raggedright
  \item \emph{Say it:} \emph{taṭṭu}
  \item \emph{Say it:} \emph{taṭṭu}, once more
  \item \emph{Say it:} say \emph{taṭṭu}
  \item \emph{Recall:} say the Telugu for to hit, then the Telugu for to cool, then say \emph{taṭṭu} again
\end{itemize}

\begin{tcolorbox}[breakable,colback=teal!4,colframe=teal!35!black,title={Before you move on}]
What does \te{తట్టు} mean? (\textquotedblleft{}To tap, to pat\textquotedblright{}.) Say it once more.
\end{tcolorbox}
\section[\texorpdfstring{\te{తెంపు} (tempu)}{తెంపు (tempu)}]{\te{తెంపు} (tempu) — to pluck, to snap off}

\label{lesson:TE-C207-tempu}



\begin{quote}
Before the new one: say the Telugu for to cool, then the Telugu for to tap.
\end{quote}

\subsection*{You'll want to know: \te{తెంపు}}
\textbf{\te{తెంపు}} — \emph{tempu} — \textquotedblleft{}to pluck, to snap off\textquotedblright{}.

\begin{grammarlens}[title={What you've built}]
One more word for this chapter.
\end{grammarlens}

\subsection*{Your turn}
\begin{itemize}
\raggedright
  \item \emph{Say it:} \emph{tempu}
  \item \emph{Say it:} \emph{tempu}, once more
  \item \emph{Say it:} say \emph{tempu}
  \item \emph{Recall:} say the Telugu for to cool, then the Telugu for to tap, then say \emph{tempu} again
\end{itemize}

\begin{tcolorbox}[breakable,colback=teal!4,colframe=teal!35!black,title={Before you move on}]
What does \te{తెంపు} mean? (\textquotedblleft{}To pluck, to snap off\textquotedblright{}.) Say it once more.
\end{tcolorbox}
\section[\texorpdfstring{\te{దున్ను} (dunnu)}{దున్ను (dunnu)}]{\te{దున్ను} (dunnu) — to till (a field)}

\label{lesson:TE-C207-dunnu}



\begin{quote}
Before the new one: say the Telugu for to tap, then the Telugu for to pluck.
\end{quote}

\subsection*{You'll want to know: \te{దున్ను}}
\textbf{\te{దున్ను}} — \emph{dunnu} — \textquotedblleft{}to till (a field)\textquotedblright{}.

\begin{grammarlens}[title={What you've built}]
One more word for this chapter.
\end{grammarlens}

\subsection*{Your turn}
\begin{itemize}
\raggedright
  \item \emph{Say it:} \emph{dunnu}
  \item \emph{Say it:} \emph{dunnu}, once more
  \item \emph{Say it:} say \emph{dunnu}
  \item \emph{Recall:} say the Telugu for to tap, then the Telugu for to pluck, then say \emph{dunnu} again
\end{itemize}

\begin{tcolorbox}[breakable,colback=teal!4,colframe=teal!35!black,title={Before you move on}]
What does \te{దున్ను} mean? (\textquotedblleft{}To till (a field)\textquotedblright{}.) Say it once more.
\end{tcolorbox}
\section[\texorpdfstring{\te{పండించు} (paṇḍiñcu)}{పండించు (paṇḍiñcu)}]{\te{పండించు} (paṇḍiñcu) — to grow (crops)}

\label{lesson:TE-C207-pandincu}



\begin{quote}
Before the new one: say the Telugu for to pluck, then the Telugu for to till.
\end{quote}

\subsection*{You'll want to know: \te{పండించు}}
\textbf{\te{పండించు}} — \emph{paṇḍiñcu} — \textquotedblleft{}to grow (crops)\textquotedblright{}.

\begin{grammarlens}[title={What you've built}]
One more word for this chapter.
\end{grammarlens}

\subsection*{Your turn}
\begin{itemize}
\raggedright
  \item \emph{Say it:} \emph{paṇḍiñcu}
  \item \emph{Say it:} \emph{paṇḍiñcu}, once more
  \item \emph{Say it:} say \emph{paṇḍiñcu}
  \item \emph{Recall:} say the Telugu for to pluck, then the Telugu for to till, then say \emph{paṇḍiñcu} again
\end{itemize}

\begin{tcolorbox}[breakable,colback=teal!4,colframe=teal!35!black,title={Before you move on}]
What does \te{పండించు} mean? (\textquotedblleft{}To grow (crops)\textquotedblright{}.) Say it once more.
\end{tcolorbox}
\section[(dialogue)]{Words from chapters 199-202 — a first pass}

\label{lesson:TE-R207-first-pass-words-from-chapters-199-202}



\begin{quote}
Say the words for to till and to grow.
\end{quote}

\subsection*{Your turn}
Say these aloud:

\begin{itemize}
\raggedright
  \item to pound, to fill, to put out to dry, to fold, to arrange
  \item to save, to choose, to try, to serve, to chew
  \item to be found, to lie down, to yawn, to shiver, to move
  \item to discuss, to contact, to check, to register, to practise
\end{itemize}

\begin{tcolorbox}[breakable,colback=teal!4,colframe=teal!35!black,title={Before you move on}]
To pound? (\textbf{\te{దంచు}}, \emph{dañcu}.) To fill? (\textbf{\te{నింపు}}, \emph{nimpu}.) To put out to dry? (\textbf{\te{ఆరబెట్టు}}, \emph{ārabeṭṭu}.) To fold? (\textbf{\te{మడతపెట్టు}}, \emph{maḍatapeṭṭu}.) To arrange? (\textbf{\te{సర్దు}}, \emph{sardu}.) To save? (\textbf{\te{పొదుపు} \te{చేయు}}, \emph{podupu cēyu}.) To choose? (\textbf{\te{ఎంచుకో}}, \emph{eñcukō}.) To try? (\textbf{\te{ప్రయత్నించు}}, \emph{prayatniñcu}.) To serve? (\textbf{\te{వడ్డించు}}, \emph{vaḍḍiñcu}.) To chew? (\textbf{\te{నమలు}}, \emph{namalu}.) To be found? (\textbf{\te{దొరుకు}}, \emph{doruku}.) To lie down? (\textbf{\te{పడుకో}}, \emph{paḍukō}.) To yawn? (\textbf{\te{ఆవలించు}}, \emph{āvaliñcu}.) To shiver? (\textbf{\te{వణుకు}}, \emph{vaṇuku}.) To move? (\textbf{\te{కదులు}}, \emph{kadulu}.) To discuss? (\textbf{\te{చర్చించు}}, \emph{carciñcu}.) To contact? (\textbf{\te{సంప్రదించు}}, \emph{sampradiñcu}.) To check? (\textbf{\te{తనిఖీ} \te{చేయు}}, \emph{tanikhī cēyu}.) To register? (\textbf{\te{నమోదు} \te{చేయు}}, \emph{namōdu cēyu}.) To practise? (\textbf{\te{సాధన} \te{చేయు}}, \emph{sādhana cēyu}.)
\end{tcolorbox}
\section[(dialogue)]{Words from chapters 203-207 — a second pass}

\label{lesson:TE-R207-second-pass-words-from-chapters-203-207}



\begin{quote}
Say the words for to scold and to fight.
\end{quote}

\subsection*{Your turn}
Say these aloud:

\begin{itemize}
\raggedright
  \item to scold, to fight, to grow up, to bloom, to rain down
  \item to compare, to fall, to exercise, to act, to draw
  \item to record, to take a photo, to share, to find out, to borrow
  \item to roll, to slip, to hang, to tread on, to hit
  \item to cool, to tap, to pluck, to till, to grow
\end{itemize}

\begin{tcolorbox}[breakable,colback=teal!4,colframe=teal!35!black,title={Before you move on}]
To scold? (\textbf{\te{తిట్టు}}, \emph{tiṭṭu}.) To fight? (\textbf{\te{పోట్లాడు}}, \emph{pōṭlāḍu}.) To grow up? (\textbf{\te{ఎదుగు}}, \emph{edugu}.) To bloom? (\textbf{\te{పూయు}}, \emph{pūyu}.) To rain down? (\textbf{\te{కురియు}}, \emph{kuriyu}.) To compare? (\textbf{\te{పోల్చు}}, \emph{pōlcu}.) To fall? (\textbf{\te{రాలు}}, \emph{rālu}.) To exercise? (\textbf{\te{వ్యాయామం} \te{చేయు}}, \emph{vyāyāmaṁ cēyu}.) To act? (\textbf{\te{నటించు}}, \emph{naṭiñcu}.) To draw? (\textbf{\te{గీయు}}, \emph{gīyu}.) To record? (\textbf{\te{రికార్డు} \te{చేయు}}, \emph{rikārḍu cēyu}.) To take a photo? (\textbf{\te{ఫోటో} \te{తీయు}}, \emph{phōṭō tīyu}.) To share? (\textbf{\te{పంచుకో}}, \emph{pañcukō}.) To find out? (\textbf{\te{కనుక్కో}}, \emph{kanukkō}.) To borrow? (\textbf{\te{అప్పు} \te{తీసుకో}}, \emph{appu tīsukō}.) To roll? (\textbf{\te{దొర్లు}}, \emph{dorlu}.) To slip? (\textbf{\te{జారు}}, \emph{jāru}.) To hang? (\textbf{\te{వేలాడు}}, \emph{vēlāḍu}.) To tread on? (\textbf{\te{తొక్కు}}, \emph{tokku}.) To hit? (\textbf{\te{కొట్టు}}, \emph{koṭṭu}.) To cool? (\textbf{\te{చల్లార్చు}}, \emph{callārcu}.) To tap? (\textbf{\te{తట్టు}}, \emph{taṭṭu}.) To pluck? (\textbf{\te{తెంపు}}, \emph{tempu}.) To till? (\textbf{\te{దున్ను}}, \emph{dunnu}.) To grow? (\textbf{\te{పండించు}}, \emph{paṇḍiñcu}.)
\end{tcolorbox}
